\hypertarget{the-remaining-surfaces}{%
\chapter{12. The remaining surfaces}\label{the-remaining-surfaces}}

The resource chapters covered the large outward-facing classes, the
things a workload reaches for. What is left is smaller and inward: the
surfaces that are properties of the box the workload runs in rather than
capabilities it asks for, each one minor next to the filesystem or the
network and each one load-bearing in the threat model. They are what the
workload can observe of its neighbours, what it inherits at startup, how
its own terminal can be turned against the operator, and what it might
use to climb or to map the machine.

A note on what is settled here. These surfaces are shipped and enforced
as described, with one class of exception: a few convenience toggles
over them are not built yet, and the parser rejects them today. A file
to seed the environment from, a template to seed the home's dotfiles,
and named shortcuts for GPU and hardware-token access are all roadmap.
The controls they would name are reachable now through the general
primitives, a forced environment variable, a device passthrough, a
socket grant, so nothing in the threat model waits on them, and the
toggles are sugar over mechanisms that already hold. Where the text
reaches one of them, it says so.

\hypertarget{what-the-workload-can-observe}{%
\section{12.1 What the workload can
observe}\label{what-the-workload-can-observe}}

A workload that can look into other processes need not reach the network
to do harm. One that can \texttt{ptrace} the operator's shell reads its
memory directly, every password and decrypted key and session token in
it; one that can read another process's
\texttt{/proc/\textless{}pid\textgreater{}/environ} lifts whatever
secret was passed to it that way; one that can signal arbitrary
processes can kill the daemons that were watching it (T1.1). The
controls against all three are structural rather than declared. The PID
namespace is the strong one: a process the kennel cannot see in
\texttt{/proc} it cannot name to \texttt{ptrace} or to \texttt{kill},
and \texttt{/proc} is mounted so that even visible entries reveal
nothing of their owners. Signal delivery is fenced a second way on a
recent kernel, by a Landlock scope that refuses signalling any process
outside the sandbox even when a PID is guessed, the same scoping family
that closed the abstract-socket namespace two chapters back.

The policy carries \texttt{{[}unsafe.ptrace{]}} and
\texttt{{[}unsafe.signal{]}} tables, and they are worth being clear
about, because they are advisory. They state which targets a workload
may reach, and that statement is real in that the diff shows it and the
operator signs it, but the enforcement underneath is the PID namespace
and seccomp, not the tables. Declaring one compiles with a warning
(\texttt{footgun-warn-dont-forbid}): it records an intent, it does not
impose a control the structure was not already imposing. The distinction
matters wherever these surfaces appear, because several of them are held
by construction and merely named by policy, and reading the policy as
the enforcement would overstate what a table can do.

\hypertarget{what-the-workload-inherits}{%
\section{12.2 What the workload
inherits}\label{what-the-workload-inherits}}

The environment a process starts with is a quiet credential surface,
because the operator's shell is usually full of secrets,
\texttt{AWS\_SECRET\_ACCESS\_KEY}, \texttt{GITHUB\_TOKEN}, an agent
socket path, and a child inherits all of them by default. Kennel does
not inherit and filter, it synthesises. The spawn builds the workload's
environment from policy and execs with that, never consulting the
parent's environment as a source of truth, so the default is empty and
every variable present is one policy named, set in
\texttt{{[}env{]}.set}. This is the denylist lesson applied to the
environment: curating the dangerous variables out of the parent's set
leaks the one the list forgot, while building the set from nothing leaks
only what was written down (\texttt{construction-by-absence}). A
workload that truly needs a dynamic non-secret like \texttt{TERM} opts
that single variable in, visibly, as a discouraged exception (T1.1).

The run context around the environment follows the same rule.
\texttt{\$PATH} is set from the execute allowlist rather than inherited,
because it means nothing except as an index into the binaries that
allowlist already permits. The login shell is policy-selectable but must
be a binary the allowlist permits, so naming a shell the kennel could
not then run is a compile error and not a runtime surprise
(\texttt{make-invalid-unrepresentable}). The shell-init files are
synthesised from policy and rebuilt every spawn, in a home that is the
kennel's own and never the operator's. That rebuild is a security
property rather than a convenience: a persistent, workload-writable
\texttt{\textasciitilde{}/.bashrc} runs on every future interactive
shell in the kennel, which makes it a self-persistence vector, so the
home is reconstructed by default and a policy that wants durable state
names the exact paths that survive. Persisting a working directory is
the ordinary low-risk case; persisting a shell-init file is a
deliberate, named, diff-reviewed choice, and even then the blast radius
is this kennel's own later runs, re-entered under the same confinement.

\hypertarget{the-terminal}{%
\section{12.3 The terminal}\label{the-terminal}}

A workload running in a terminal can turn that terminal against the
operator two ways, and they need different answers. The first is input
injection. The \texttt{TIOCSTI} ioctl asks the kernel to push characters
into the controlling terminal as though typed, which a confined process
would use to queue commands that run in the operator's shell after the
kennel exits. Kennel closes it by removing the target rather than
policing the call: the workload's controlling terminal is a node in the
kennel's own isolated \texttt{devpts}, and the operator's
\texttt{/dev/tty} is not in the view at all, so an injection reaches
only the workload's own session, which is the workload's to disrupt
(\texttt{construction-by-absence}). No syscall filter is needed, because
there is nothing on the other side to inject into.

The second is output. A workload writing to its terminal can emit escape
sequences that act on the operator's emulator rather than printing, OSC
52 to read and write the system clipboard, OSC 9 and 777 to raise
notifications, the device-control bands to carry emulator-specific
commands, all of it an exfiltration and injection channel with nothing
to do with \texttt{TIOCSTI}. The output is filtered on its way to the
operator's real terminal, the dangerous set dropped and the benign
passed, and where that filter runs is the careful part. The daemon's PTY
broker pumps the workload's output as raw bytes and never parses it, so
the parser that interprets a hostile escape stream, a streaming
\texttt{vte} state machine, lives in the attached client and not in the
trusted base (\texttt{quarantine-the-unsafe}); the daemon conveys the
decision and the client enforces it. The workload cannot choose its
client, so it cannot opt out, and a consumer that bypasses the filter
only exposes its own terminal. The filter is best-effort and not a
proof: it shuts the low-effort clipboard and notification channel (T2.6)
rather than every conceivable desync, and running a kennel in a terminal
you trust stays the backstop.

The controlling terminal an interactive run needs is built to preserve
that isolation. A real terminal is required for job control and
full-screen tools, but forwarding the operator's own would hand the
workload the very \texttt{/dev/tty} the input defence removes. So the
pty is allocated from the kennel's own \texttt{devpts} after the pivot,
the seal makes it the workload's controlling terminal and hands the
master not to the CLI but to the daemon, which holds it in a per-kennel
broker and pumps it raw to whatever client is attached. Because the
master lives in the daemon, the operator can detach and reattach without
ending the workload, the console model rather than an exec into the live
sandbox: no second process is injected past the confinement, only a
terminal taken over and handed back.

\hypertarget{capabilities-mounts-and-the-floor}{%
\section{12.4 Capabilities, mounts, and the
floor}\label{capabilities-mounts-and-the-floor}}

The last surfaces are the ones a workload might use to climb or to map
the machine, and the answers are uniform and dull by design. Linux
capabilities are dropped to nothing: the bounding set is emptied so a
setuid binary cannot raise the workload past its uid,
\texttt{no\_new\_privs} is forced on as the execution chapter required,
and the schema warns at the mere listing of a capability, because a
confined uid-1000 workload that needs one is almost always a design
error (\texttt{split-the-uid}). The mount table is narrowed by the same
mount namespace the constructed view already uses: the workload sees
only the mounts it was given, none of the operator's removable media or
cloud sync or loop images, and a \texttt{mount} call fails for want of
the capability just dropped.

Beneath those sits the syscall floor. Most of what Kennel confines reads
better as a resource ACL than as a syscall filter, so seccomp is not the
main event, but it closes the gaps the ACLs cannot reach: the
abstract-socket connect Landlock has no word for, and a denylist of
calls with no legitimate use in a kennel and a long history in exploit
chains, \texttt{userfaultfd}, \texttt{bpf}, \texttt{ptrace}, the
\texttt{process\_vm} pair, \texttt{mount} and \texttt{pivot\_root} and
the rest. The cgroup the kennel runs in is one it cannot leave, which
matters because leaving it would shed the BPF filters attached to it, so
membership is an invariant the operator does not set and the workload
cannot change. None of these is a knob the policy author reaches for;
they are the floor the surfaces above stand on.
